% ======================================================================
%  PARTIE 3 -- YOLO                                          (~1 min 30)
% ======================================================================
\section{YOLO : segmentation sémantique}

% ----------------------------------------------------------------------
% Schéma d'une scène simplifiée + points LiDAR, selon la tâche YOLO :
%   \scenesegm{1} = détection (boîtes), \scenesegm{2} = segmentation
%   d'instance, \scenesegm{3} = segmentation sémantique
% (\shorthandoff : les caractères ? : ! ; du français perturbent les calculs TikZ)
% ----------------------------------------------------------------------
\shorthandoff{;:!?}
\newcommand{\scenesegm}[1]{%
  \begin{tikzpicture}[x=0.95cm, y=0.95cm]
    \ifcase#1\or
      \def\sS{gray!8}\def\sB{gray!30}\def\sV{gray!22}\def\sR{gray!45}\def\sC{gray!62}%
    \or
      \def\sS{gray!8}\def\sB{gray!30}\def\sV{gray!22}\def\sR{gray!45}\def\sC{cVoiture!55}%
    \or
      \def\sS{cCiel!30}\def\sB{cBati!40}\def\sV{cVege!45}\def\sR{cRoute!45}\def\sC{cVoiture!55}%
    \fi
    \begin{scope}
      \clip (0,0) rectangle (4,2.4);
      \fill[\sS] (0,0) rectangle (4,2.4);                 % ciel
      \fill[\sB] (0,0.85) rectangle (1.5,2.0);            % bâtiment
      \fill[\sV] (3.3,1.35) ellipse (0.8 and 0.6);        % végétation
      \fill[\sR] (0,0) rectangle (4,0.85);                % route
      \fill[\sC, rounded corners=2pt] (1.6,0.25) rectangle (2.7,0.85);   % voiture
      \fill[\sC, rounded corners=2pt] (1.85,0.85) rectangle (2.45,1.15);
      \fill[black!70] (1.85,0.27) circle (0.12) (2.45,0.27) circle (0.12);
    \end{scope}
    \draw[gray!70] (0,0) rectangle (4,2.4);
    % Points LiDAR (pas de retour dans le ciel)
    \foreach \y in {0.12,0.37,0.62,0.87,1.12,1.37,1.62,1.87} {
      \foreach \x in {0.14,0.42,0.70,0.98,1.26,1.54,1.82,2.10,2.38,2.66,2.94,3.22,3.50,3.78} {
        % région du point : 1 voiture, 2 végétation, 3 bâtiment, 4 route, 0 ciel
        % (produits de booléens plutôt que && : la figure est dans un tableau)
        \pgfmathtruncatemacro{\sCar}{max((\x>=1.6)*(\x<=2.7)*(\y>=0.25)*(\y<=0.85),
                                         (\x>=1.85)*(\x<=2.45)*(\y>0.85)*(\y<=1.15))}%
        \pgfmathtruncatemacro{\sVeg}{((\x-3.3)/0.8)^2+((\y-1.35)/0.6)^2<=1}%
        \pgfmathtruncatemacro{\sBat}{(\x<=1.5)*(\y>0.85)*(\y<=2.0)}%
        \pgfmathtruncatemacro{\sRou}{\y<=0.85}%
        \pgfmathtruncatemacro{\sBoite}{(\x>=1.4)*(\x<=2.9)*(\y>=0.05)*(\y<=1.3)}%
        \def\sReg{0}%
        \ifnum\sRou=1 \def\sReg{4}\fi
        \ifnum\sBat=1 \def\sReg{3}\fi
        \ifnum\sVeg=1 \def\sReg{2}\fi
        \ifnum\sCar=1 \def\sReg{1}\fi
        \ifnum\sReg>0
          \def\sCoul{none}%
          \ifcase#1\or
            \ifnum\sBoite=1 \def\sCoul{cVoiture}\fi
          \or
            \ifnum\sReg=1 \def\sCoul{cVoiture}\fi
          \or
            \ifcase\sReg\or\def\sCoul{cVoiture}\or\def\sCoul{cVege}\or\def\sCoul{cBati}\or\def\sCoul{cRoute}\fi
          \fi
          \def\sAucune{none}%
          \ifx\sCoul\sAucune
            \filldraw[fill=white, draw=gray, line width=0.3pt] (\x,\y) circle (1.4pt);
          \else
            \filldraw[fill=\sCoul, draw=white, line width=0.3pt] (\x,\y) circle (1.6pt);
          \fi
        \fi
      }
    }
    \ifnum#1=1
      \draw[accent, very thick] (1.4,0.05) rectangle (2.9,1.3);
      \node[fill=accent, text=white, font=\tiny\bfseries, anchor=south west, inner sep=1.5pt] at (1.4,1.3) {car};
    \fi
  \end{tikzpicture}}
\shorthandon{;:!?}

% ----------------------------------------------------------------------
% À DIRE (~40 s)
%  - YOLO = You Only Look Once : réseau de neurones qui analyse l'image en
%    UNE SEULE passe, contrairement aux approches plus anciennes en deux
%    étapes (ex. famille R-CNN) qui repèrent d'abord des zones d'intérêt
%    puis les analysent une par une.
%  - Beaucoup plus rapide -> choix classique pour du temps réel embarqué
%    sur robot.
%  - Plusieurs versions (YOLOv5, v8, 11, 26), chacune plus précise et plus
%    rapide que la précédente ; et plusieurs tailles de modèle, de nano
%    (la plus légère, la moins précise) à extra-large (la plus lourde, la
%    plus précise).
%  - Choix : nano -> faible coût de calcul, car la chaîne de cartographie
%    (LIO) est déjà gourmande. Point de départ volontairement
%    conservateur : on pourra passer à small / medium selon la puissance
%    réellement disponible sur le robot.
% ----------------------------------------------------------------------
\begin{frame}{YOLO : principe}
  \begin{columns}[c]
    \begin{column}{0.36\textwidth}
      \begin{itemize}
        \item \alert{Y}ou \alert{O}nly \alert{L}ook \alert{O}nce
        \item Réseau de neurones :\\ \alert{une seule passe}
        \item \alert{Rapide} : temps réel,\\ embarqué sur robot
        \item Versions : v5, v8, 11, \alert{26}
      \end{itemize}
    \end{column}
    \begin{column}{0.62\textwidth}
      \centering
      \begin{tikzpicture}[font=\scriptsize]
        % Deux étapes
        \node[anchor=west, font=\scriptsize\bfseries, text=textgray] at (-0.6,1.85) {Approche en deux étapes};
        \node[capteur=1.0cm] (i1) at (0,1.2) {Image};
        \node[boite=1.55cm] (r1) at (1.9,1.2) {Zones d'intérêt};
        \node[boite=1.55cm] (c1) at (3.95,1.2) {Analyse une\\ par une};
        \node[boite=1.05cm] (o1) at (5.75,1.2) {Résultat};
        \draw[fleche] (i1) -- (r1);
        \draw[fleche] (r1) -- (c1);
        \draw[fleche] (c1) -- (o1);
        \node[right=1pt of o1, text=kored] {\ko\ lent};
        % YOLO
        \node[anchor=west, font=\scriptsize\bfseries, text=primaryblue] at (-0.6,0.35) {YOLO};
        \node[capteur=1.0cm] (i2) at (0,-0.3) {Image};
        \node[donnee=3.6cm] (c2) at (2.925,-0.3) {\textbf{un seul réseau, une seule passe}};
        \node[boite=1.05cm] (o2) at (5.75,-0.3) {Résultat};
        \draw[fleche] (i2) -- (c2);
        \draw[fleche] (c2) -- (o2);
        \node[right=1pt of o2, text=okgreen] {\ok\ rapide};
        % Tailles de modèle
        \node[anchor=west, font=\scriptsize\bfseries, text=textgray] at (-0.6,-1.25) {Tailles de modèle};
        \draw[thick, primaryblue!50] (0.2,-2.0) -- (6.2,-2.0);
        \foreach \px/\nom in {0.4/nano,1.8/small,3.2/medium,4.6/large,6.0/x-large} {
          \fill[primaryblue] (\px,-2.0) circle (2.2pt);
          \node[below=3pt, font=\scriptsize] at (\px,-2.0) {\nom};
        }
        \draw[accent, very thick] (0.4,-2.0) circle (5pt);
        \node[above=4pt, text=accent, font=\scriptsize\bfseries] at (0.4,-2.0) {choisi};
        \node[font=\tiny, text=textgray, anchor=north] at (0.4,-2.45) {rapide, léger};
        \node[font=\tiny, text=textgray, anchor=north] at (6.0,-2.45) {précis, lourd};
      \end{tikzpicture}
    \end{column}
  \end{columns}
\end{frame}

% ----------------------------------------------------------------------
% À DIRE (~50 s)
%  - Selon la tâche, YOLO donne :
%      * une boîte autour de chaque objet (détection) ;
%      * un masque de pixels par objet détecté (segmentation d'instance) ;
%      * un label pour CHAQUE pixel, fond compris (route, végétation,
%        ciel...) : segmentation sémantique (disponible avec YOLO26).
%  - Mon besoin : donner un label à chaque point LiDAR projeté dans
%    l'image.
%      * Avec une boîte : un point dans le coin de la boîte appartient en
%        fait au fond, mais il reçoit le label de l'objet -> erreurs.
%      * L'instance corrige ce problème pour les objets, mais ne masque
%        que les objets reconnus (véhicules, piétons...) : route,
%        végétation, sol restent sans label -> la plupart des points
%        n'auraient pas de label.
%      * La sémantique étiquette tous les pixels -> tous les points.
%  - Modèle retenu : YOLO26 nano en segmentation sémantique, entraîné sur
%    Cityscapes (scènes urbaines, 19 classes : route, trottoir, bâtiment,
%    végétation, voiture, piéton...).
% ----------------------------------------------------------------------
\begin{frame}{YOLO : pourquoi la segmentation \emph{sémantique} ?}
  \centering
  \begin{tabular}{@{}c@{\hspace{0.45cm}}c@{\hspace{0.45cm}}c@{}}
    \scenesegm{1} & \scenesegm{2} & \scenesegm{3} \\[2pt]
    \textbf{Détection} (boîtes) & \textbf{Segmentation d'instance} & \textbf{Segmentation sémantique} \\
    {\footnotesize \ko\ fond étiqueté « voiture »} &
    {\footnotesize \ko\ fond non étiqueté} &
    {\footnotesize \ok\ un label par pixel} \\
  \end{tabular}\par
  \vspace{0.25cm}
  \begin{tikzpicture}[font=\scriptsize]
    \filldraw[fill=cVoiture, draw=white] (0,0) circle (1.6pt); \node[right] at (0.05,0) {point LiDAR étiqueté};
    \filldraw[fill=white, draw=gray] (3.1,0) circle (1.4pt); \node[right] at (3.15,0) {point sans label};
  \end{tikzpicture}\par
  \vspace{0.25cm}
  \begin{tcolorbox}[question, width=0.78\textwidth]
    Modèle retenu : \alert{YOLO26 nano}, segmentation sémantique,
    entraîné sur \alert{Cityscapes} {\small(scènes urbaines, 19 classes)}
  \end{tcolorbox}
\end{frame}
